\section{Cargador de arranque}
\subsection{GRUB}
\begin{frame}

    Configuramos GRUB, el cargador de arranque. Debemos tener en cuenta si la placa tiene BIOS o UEFI como vimos al principio.
    
    \only<1>{\begin{tcolorbox}[title=Si se trata de un arranque BIOS con disco mecánico o una unidad SSD]%
        \orden{grub-install /dev/sda}%
    \end{tcolorbox}}
    \only<2>{\begin{tcolorbox}[title=Si se trata de un arranque BIOS con una unidad M.2]%
        \orden{grub-install /dev/nvme0n1}%
    \end{tcolorbox}\vspace{5mm}}
    \only<3>{\begin{tcolorbox}[title=Si se trata de un arranque UEFI]
      \orden{grub-install --target=x86\_64-efi --efi-directory=/boot --bootloader-id=grub}
    \end{tcolorbox}\vspace{5mm}}

    \only<4>{\begin{tcolorbox}[title=Guardamos la configuración:]
      \orden{grub-mkconfig -o /boot/grub/grub.cfg}
      \tcblower
      \scriptsize
      Y si no muestra un mensaje de error, podemos reiniciar comprobando que podemos arrancar Arch Linux.
    \end{tcolorbox}}\pause\vspace{5mm}
     
    
\end{frame}
